\section{Conclusion}

This paper isolates activation choice as a controlled design variable in Informer forecasting. Under a matched official-script-style protocol, activation-only changes alter the best configuration in most dataset-horizon cells: non-GELU activations are best in 34 of 36 completed cells, and bounded or tanh-family configurations have the strongest static ranks. The result is practically useful because it does not require changing attention, width, depth, or the training schedule.

The evidence also sets clear limits. \texttt{tanh\_sin001}, \texttt{tanh}, and \texttt{softsign} are strong Informer candidates, but none is universal. Output \texttt{tanh} has several high-gain wins and a poor average rank. Decoder-only \texttt{tanh} is weak. Solar gains are smaller and site-specific, with Solar3 remaining GELU-best at both tested horizons. Factor diagnostics reveal regime dependence, but current leave-one-seed-out evidence does not support a deployable zero-shot activation router.

The strongest current design lesson is therefore conservative: GELU should not be treated as a harmless default in Informer-style forecasting, but activation choice should be validated per dataset, horizon, placement, and architecture. Future work should turn the diagnostic factor evidence into prospective routing experiments, expand the iTransformer placement probe to more seeds and cells, and test whether architecture-aware rules can select candidate activations before training.
